% Section 6.4, from lectures/L06/appendix/c_exercises.md. The worked solutions are not
% printed; the aside says where they are.

\section{Exercises}
\label{c6:sec:exercises}

Eight, ending with the Cross-check. That one has four legs, and one of them is a tempting argument you are asked to break.

\begin{aside}
\textbf{How to check your work.} A \kindtag{Code} exercise is judged by the suite that ships with it: run \code{make test} at the root of the companion repository, with \code{AEL_DIR} pointing at your toolkit. A suite you have not started is reported as \code{SKIP} rather than as a failure, so the command is worth running from the first day. The hand calculations and designs are checked against the toolkit functions each one names.

\textbf{The solutions are not in this book.} They are published in full in the companion repository, in \file{lectures/L06/appendix}, including the plausible wrong answer where there is one. Read them once you have your own answers, including the ones you are unsure about, because being unsure and right is a different thing from being unsure and wrong and the solutions say which is which.
\end{aside}

% ------------------------------------------------------------------
\recall{The operating point}
\label{c6:ex:1}
\begin{enumerate}
  \item Name the three numbers that describe a quiescent point, and say why it is usually two.
  \item Write the load line of a stage with both a collector and an emitter resistor.
  \item Where on it does an amplifier sit, and what decides ``where''?
  \item A stage is biased so that the collector sits at 0.5 V above the emitter. What is wrong with it, and what would the output look like?
\end{enumerate}

% ------------------------------------------------------------------
\recall{Stiffness}
\label{c6:ex:2}
\begin{enumerate}
  \item Define divider stiffness, and give the usual rule of thumb.
  \item What error in the collector current does that rule of thumb accept?
  \item What does buying a factor of ten more stiffness cost?
  \item Name two ways of making the base current negligible instead, and say which chapters use them.
\end{enumerate}

% ------------------------------------------------------------------
\handcalc{The quiescent point, twice}
\label{c6:ex:3}
For the stage in \secref{c6:sec:a2}: 10 V, 33 kilohm over 6.8 kilohm, 1 kilohm emitter, 4.7 kilohm collector, $\beta = 50$.

\begin{enumerate}
  \item The base voltage, emitter voltage and collector current, ignoring base current.
  \item The base current that implies, and the droop it causes across the divider's Thevenin resistance.
  \item The corrected base voltage and collector current. One iteration is enough; say why.
  \item The percentage error the first answer made, and the stiffness of this divider.
\end{enumerate}
\checkyourself{\code{quiescentPoint}, \code{stiffness}}

% ------------------------------------------------------------------
\handcalc{Breaking the tempting argument}
\label{c6:ex:4}
The argument in \secref{c6:sec:b2} claims a 10 per cent current rise drops $V_{BE}$ from 0.65 V to 0.55 V, which then restores the current.

\begin{enumerate}
  \item What collector current does $V_{BE} = 0.55$ V actually give, relative to 0.65 V?
  \item Is that consistent with ``the current rose 10 per cent''?
  \item Identify the step in the argument that is wrong. It is not the 10 per cent.
  \item Write the correct version in two sentences, using the word ``loop''.
\end{enumerate}
\checkyourself{\code{currents}}

% ------------------------------------------------------------------
\design{Bias a stage to a specification}
\label{c6:ex:5}
Design a common-emitter stage: 12 V supply, 2 mA quiescent collector current, the collector sitting at half the supply, and drift below 0.3 per cent per degree.

\begin{enumerate}
  \item Choose the collector resistor from the quiescent voltage requirement.
  \item Choose the emitter resistor from the drift requirement, then check it against the 220 mV rule.
  \item Choose the divider for a stiffness of at least 20, in E12 values.
  \item State the quiescent point your design actually achieves, including the base-current droop.
  \item State the emitter factor, and therefore what this design has given up in gain.
\end{enumerate}
\checkyourself{\code{quiescentPoint}, \code{stiffness}, \code{degenerationResistor}, \code{emitter_factor}}

% ------------------------------------------------------------------
\codeexercise{The bias point}
\label{c6:ex:6}
Implement \code{ael::bias} to the specification in \secref{c6:sec:b5}.

\code{quiescentPoint} is the only one that needs thought, because it is self-referential: the base current depends on the collector current which depends on the base voltage which depends on the base current. Iterate it. Two or three passes converge to well under a per cent.

% ------------------------------------------------------------------
\codeexercise{Temperature in the device}
\label{c6:ex:7}
Add \code{temperature} to \code{ael::device::bjt::Parameters} and make both the thermal voltage and the saturation current depend on it.

Then check the model against the number everyone quotes: at 1 mA and fixed collector current, the base-emitter voltage should fall by about 1.8 mV per degree. \textbf{If yours gives exactly 2.0, you have put the coefficient in by hand rather than letting it emerge}, and the Cross-check below will not work.

% ------------------------------------------------------------------
\crosscheck{The drift, four ways}
\label{c6:ex:8}
The stage of \secref{c6:sec:a2}. How much does its collector current change per degree?

\begin{enumerate}
  \item \textbf{By the tempting argument.} Follow \secref{c6:sec:b2} to its conclusion and state what it predicts. You will not be able to.
  \item \textbf{By hand, linearised.} The base-emitter voltage drifts $-2$ mV per degree at a held base, so the emitter voltage rises 2 mV and the current rises by 2 mV over the emitter resistor.
  \item \textbf{By your solver.} Sweep the temperature from 27 to 37 degrees and measure.
  \item \textbf{With the emitter resistor removed}, by the solver, for comparison.
\end{enumerate}

\task{What to expect}
\textbf{Leg 1 does not produce a number.} Followed honestly it says the current falls by a factor of 47, which contradicts the 10 per cent rise it started from. That is the exercise: an argument can be qualitatively right, easy to believe, and still not close.

\textbf{Leg 2 gives about 0.21 per cent per degree.}

\textbf{Leg 3 gives about 0.17 per cent per degree}, roughly 20 to 30 per cent below leg 2.

\textbf{That gap is real and has two named causes}, and neither is an error:

\begin{itemize}
  \item The $-2$ mV per degree is a round figure. The physics gives $-1.77$ mV per degree at 1 mA, and the coefficient itself depends on the operating current.
  \item Leg 2 holds the base voltage fixed. The solver lets it droop further as the current rises, because \secref{c6:sec:a3}'s base current rises with it. That is a second loop the hand calculation does not contain.
\end{itemize}
\textbf{Leg 4 gives about 7 per cent per degree}, forty times worse, and that is the number the whole chapter exists to prevent.

\textbf{The ratio between legs 3 and 4 should be about the emitter factor}, and checking that is the real content of the Cross-check. If it is not, the suppression is coming from somewhere other than where the theory says.

\textbf{If leg 3 gives exactly leg 2's answer}, the temperature coefficient was put into the model by hand instead of emerging from $I_S(T)$ and $V_T(T)$. That is worth catching: a model that agrees with the hand calculation exactly is a model that contains the hand calculation.
